\sect{आषाढ-कृष्ण-योगिनी-एकादशी-माहात्म्यम्}
\label{sec:padma-ashadha-krishna-yogini}

\ganapatyadiDhyanam
\padmaPuranam

\dnsub{अथ चतुष्पञ्चाशत्तमोऽध्यायः॥५४}

\uvacha{युधिष्ठिर उवाच}

\twolineshloka
{आषाढकृष्णपक्षे तु किं वै एकादशी भवेत्}
{कथयस्व प्रसादेन वासुदेव ममाग्रतः}% ॥१॥

\uvacha{श्रीकृष्ण उवाच}

\twolineshloka
{व्रतानामुत्तमं राजन् कथयामि तवाग्रतः}
{सर्वपापक्षयकरं सर्वमुक्तिप्रदायकम्}% ॥२॥

\twolineshloka
{आषाढस्यासिते पक्षे योगिनी नाम नामतः}
{एकादशी नृपश्रेष्ठ महापातकनाशिनी}% ॥३॥

\twolineshloka
{संसारार्णवमग्नानां पोतभूता सनातनी}
{जगत्त्रये सारभूता योगिनी व्रतकारिणाम्}% ॥४॥

\twolineshloka
{कथयामि तवाग्रेऽहं कथां पौराणिकीं शुभाम्}
{अलकायां राजराजः शिवभक्तिपरायणः}% ॥५॥

\twolineshloka
{तस्याऽऽसीत् पुष्पबटुको हेममालीति नामतः}
{तस्य पत्नी सुरूपा च विशालाक्षीति नामतः}% ॥६॥

\twolineshloka
{स तस्यां चासक्तमना कामपाशवशं गतः}
{मानसात् पुष्पनिचयमानीय स्वगृहे स्थितः}% ॥७॥

\twolineshloka
{पत्नीप्रेमरसासक्तो न कुबेरालयं गतः}
{कुबेरो देवसदने करोति शिवपूजनम्}% ॥८॥

\twolineshloka
{मध्याह्नसमये राजन् पुष्पागमसमीक्षकः}
{हेममाली स्वभवने रमते कान्तया सह}% ॥९॥

\threelineshloka
{यक्षराट् प्रत्युवाचाथ कालातिक्रमकोपितः}
{कस्मान्नाऽऽयाति भो यक्षा हेममाली दुरात्मवान्}
{निश्चयः क्रियतामस्य इत्युवाच पुनः पुनः}% ॥१०॥

\uvacha{यक्षा ऊचुः}

\twolineshloka
{वनिताकामुको गेहे रमते स्वेच्छया नृप}
{तेषां वाक्यं समाकर्ण्य कुबेरः कोपपूरितः}% ॥११॥

\twolineshloka
{आह्वयामास तं तूर्णं बटुकं हेममालिनम्}
{ज्ञात्वा कालात्ययं सोऽपि भयव्याकुललोचनः}% ॥१२॥

\twolineshloka
{अस्नात एव आगत्य कुबेरस्याग्रतः स्थितः}
{तं दृष्ट्वा धनदः क्रुद्धः क्रोधसंरक्तलोचनः}% ॥१३॥
\onelineshloka*
{प्रत्युवाच रुषाविष्टः कोपप्रस्फुरिताधरः}

\uvacha{धनद उवाच}
\onelineshloka
{आः पाप दुष्ट दुर्वृत्त कृतवान् देवहेलनम्}% ॥१४॥

\twolineshloka
{अष्टादशकुष्ठवृतो वियुक्तः कान्तया तया}
{अस्मात् स्थानादपध्वस्तो गच्छस्व प्रमथाधम}% ॥१५॥

\twolineshloka
{इत्युक्ते वचने तस्य तस्मात् स्थानात् पपात सः}
{महादुःखाभिभूतश्च कुष्ठैः पीडितविग्रहः}% ॥१६॥

\twolineshloka
{न सुखं दिवसे तस्य न निद्रां लभते निशि}
{छायायां पीडिततनुर्निदाघेऽत्यन्तपीडितः}% ॥१७॥

\twolineshloka
{शिवपूजाप्रभावेन स्मृतिस्तस्य न लुप्यते}
{पातकेनाभिभूतोऽपि पूर्वं कर्म स्मरत्यसौ}% ॥१८॥

\twolineshloka
{भ्रममाणस्ततो गच्छन् हिमाद्रिं पर्वतोत्तमम्}
{तत्रापश्यन्मुनिवरं मार्कण्डेयं तपोनिधिम्}% ॥१९॥

\twolineshloka
{यस्याऽऽयुर्विद्यते राजन् ब्रह्मणो वयसा समम्}
{ववन्दे चरणौ तस्य दूरतः पापकर्मकृत्}% ॥२०॥

\twolineshloka
{मार्कण्डेयो मुनिवरो दृष्ट्वा तं कम्पितं तथा}
{परोपकरणार्थाय समाहूयेदमब्रवीत्}% ॥२१॥

\twolineshloka
{कस्मात् कुष्ठाभिभूतस्त्वं कुतो निन्द्यतरो ह्यसि}
{इत्युक्तः स प्रत्युवाच मार्कण्डेयं महामुनिम्}% ॥२२॥

\uvacha{हेममाल्युवाच}

\twolineshloka
{राजराजस्यानुचरो हेममालीति नामतः}
{मानसात् पद्मनिचयमानीय प्रत्यहं मुने}% ॥२३॥

\twolineshloka
{शिवपूजनवेलायां कुबेराय समर्पये}
{एकस्मिन् दिवसे चैव कालश्चाविदितो मया}% ॥२४॥

\twolineshloka
{पत्नीसौख्यप्रसक्तेन शोकव्याकुलचेतसा}
{ततः क्रुद्धेन शप्तोऽस्मि राजराजेन वै मुने}% ॥२५॥

\twolineshloka
{कुष्ठाभिभूतः सञ्जातो वियुक्तः कान्तया तया}
{अधुना तव सान्निध्यं प्राप्तोऽस्मि शुभकर्मणा}% ॥२६॥

\twolineshloka
{सतां स्वभावतश्चित्तं परोपकरणे क्षमम्}
{इति ज्ञात्वा मुनिश्रेष्ठ मां प्रशाधि कृतागसम्}% ॥२७॥

\uvacha{मार्कण्डेय उवाच}

\twolineshloka
{त्वया सत्यमिह प्रोक्तं नासत्यं भाषितं यतः}
{अतो व्रतोपदेशं ते कथयामि शुभप्रदम्}% ॥२८॥

\twolineshloka
{आषाढे कृष्णपक्षे तु योगिनीव्रतमाचर}
{अस्य व्रतस्य पुण्येन कुष्ठं यास्यति वै ध्रुवम्}% ॥२९॥

\twolineshloka
{इति वाक्यमृषेः श्रुत्वा दण्डवत् पतितो भुवि}
{उत्थापितः स मुनिना बभूवातीव हर्षितः}% ॥३०॥

\twolineshloka
{मार्कण्डेयोपदेशेन व्रतं तेन कृतं यथा}
{अष्टादशेव कुष्ठानि गतानि तस्य सर्वशः}% ॥३१॥

\twolineshloka
{मुनेर्वाचा ततः सम्यग् व्रते चीर्णेऽभवत् सुखी}
{ईदृग्विधं नृपश्रेष्ठ कथितं योगिनीव्रतम्}% ॥३२॥

\twolineshloka
{अष्टाशीतिसहस्राणि द्विजान् भोजयते तु यः}
{तत् समं फलमाप्नोति योगिनीव्रतकृन्नरः}% ॥३३॥

\twolineshloka
{महापापप्रशमनं महापुण्यफलप्रदम्}
{पठनाच्छ्रवणान्मर्त्यः सर्वपापैः प्रमुच्यते}% ॥३४॥

॥इति श्रीपाद्मे महापुराणे पञ्चपञ्चाशत्साहस्र्यां संहितायामुत्तरखण्डे उमापतिनारदसंवादान्तर्गतकृष्णयुधिष्ठिरसंवादे आषाढ-कृष्ण-योगिनी-एकादशी-माहात्म्यं नाम चतुष्पञ्चाशत्तमोऽध्यायः॥५४॥

\genMangalaShloka

\hyperref[sec:ekadashi_mahatmyam_padma_puranam]{\closesub}
\clearpage

